\documentclass[10pt]{mypackage}

\usepackage[uprightgreeks,light,noDcommand]{kpfonts}
\renewcommand{\coloneq}{\coloneqq}
\renewcommand{\eqcolon}{\eqqcolon}
%\renewcommand{\mathbb}{\mathds}
%\usepackage{homework}
\usepackage{notes}

%\usepackage[ backend=bibtex, style = alphabetic, sorting=ynt ]{biblatex}
%\addbibresource{  }

\usepackage{parskip}

\fancyhf{}
\fancyhead[R]{Avinash Iyer}
\fancyhead[L]{Probability Theory: Class Notes}
\fancyfoot[C]{\thepage}

\setcounter{secnumdepth}{0}

\begin{document}
\RaggedRight
These are notes I'm taking for Christian Gromoll's Probability Theory class. The primary text is \textit{Probability with Martingales} by Williams.
\section{Dictionary between Measure Theory and Probability}%
We will assume a lot of measure theory background, so we will instead discuss the dictionary between terms from probability and terms from measure theory. It's important to note that the measure spaces in probability theory are probability measure spaces (i.e., they have measure $1$).
\begin{center}
  \renewcommand{\arraystretch}{1.5}
  \begin{tabular}{c|c}
    Measure Theory & Probability\\
    \hline
    Set $X$ & Sample space $\Omega$\\
    $\sigma$-Algebra $\Sigma$ & $\sigma$-Field $\mathcal{F}$\\
    Measurable subsets $A\in \Sigma$ & Events $F\in \mathcal{F}$\\
    Measure $\mu(A)$ & Probability $P(F)$\\
    Almost Everywhere (a.e.) & Almost Surely (a.s.)\\
    $\ds \limsup_{n\rightarrow\infty} E_n$ & $\set{E_n\text{ infinitely often}}$, $\set{E_n\text{ i.o.}}$\\
    $\ds \liminf_{n\rightarrow\infty} E_n$ & $\set{E_n\text{ eventually}}$, $ \set{E_n\text{ ev}} $\\
    Borel Measurable Function $f\colon X\rightarrow \R$ & Random Variable $X\colon \Omega\rightarrow \R$\\
    Pushforward measure $f_{\ast}\mu$ & Law $\Lambda_X$\\
    \hline
  \end{tabular}
\end{center}
\subsection{Important Theorems}%
\begin{definition}
  Let $S$ be a set, and let $\mathcal{D}$ be a collection of subsets. We call $ \mathcal{D} $ a $\pi$-system on $S$ if it is closed under finite intersection.

  We call $\mathcal{D}$ a $\lambda$-system if
  \begin{itemize}
    \item $S\in \mathcal{D}$;
    \item for any $A,B\in \mathcal{D}$ with $A\subseteq B$, one has $B\setminus A\in \mathcal{D}$;
    \item if $\left( A_n \right)_n\subseteq \mathcal{D}$ is such that $\bigcup_n A_n = A$, then $A\in \mathcal{D}$.
  \end{itemize}
  For an arbitrary collection of subsets $\mathcal{I}$, we write $\pi\left( \mathcal{I} \right)$ for the $\pi$-system generated by $ \mathcal{I} $ and $\lambda\left( \mathcal{I} \right)$ for the $\lambda$-system generated by $ \mathcal{I} $.
\end{definition}
\begin{proposition}
  A collection $\Sigma$ of subsets of $S$ is a $\sigma$-algebra if and only if $\Sigma$ is both a $\pi$-system and a $\lambda$-system.
\end{proposition}
\begin{proof}
  Let $\Sigma$ be a $\pi$-system and a $\lambda$-system. If $E,F\in \Sigma, \left( E_n \right)_n\subseteq \Sigma$, then 
  \begin{align*}
    E^{c} = S\setminus E\\
    E\cup F &= S\setminus \left( E^{c}\cap F^{c} \right).
  \end{align*}
  Therefore, $G_n = E_1\cup\cdots\cup E_n\in \Sigma$, with $\bigcup_n G_n = \bigcup_k E_k$, so that $\bigcup_{k} E_k\in \Sigma$. Finally, we have $\bigcap_{k} E_k = \left( \bigcup_{k}E_k^{c} \right)^{c}$, meaning that $\Sigma$ is a $\sigma$-algebra.
\end{proof}
\begin{theorem}[Dynkin $\pi$-$\lambda$ Theorem]
  If $ \mathcal{I} $ is a $\pi$-system, then $\lambda\left( \mathcal{I} \right) = \sigma\left( \mathcal{I} \right)$.
\end{theorem}
\begin{proof}
  To prove the theorem, we will use a classic technique by defining a ``truth set'' on which a desired statement holds, then showing that truth set is the desired $\sigma$-algebra.

  Define
  \begin{align*}
    \mathcal{D}_1 &= \set{B\in \lambda\left( \mathcal{I} \right) : B\cap C\in \lambda\left( \mathcal{I} \right)\text{  for all }C\in \mathcal{I}}.
  \end{align*}
  We observe that $\mathcal{I}$ is a $\pi$ system, so that $\mathcal{I}\subseteq \mathcal{D}_1$. Furthermore, 
  \begin{itemize}
    \item $S\in \mathcal{D}_1$;
    \item for any $B_1,B_2\in \mathcal{D}_1$ with $B_1\subseteq B_2 $, we have that for any given $C\in \mathcal{I}$, $\left( B_2\setminus B_1 \right)\cap C = \left( B_2\cap C \right)\setminus \left( B_1\cap C \right)$, so since $\lambda\left( \mathcal{I} \right)$ is a $\lambda$-system, it follows that $B_2\setminus B_1\in \mathcal{D}_1$;
  \end{itemize}
\end{proof}
\end{document}
